\bookchapter{Capítulo 63: Views, estado e navegação}

\chapterintro{Você já sabe fazer brinquedos de uma tela só. Este capítulo resolve duas questões de “crescimento”: como pular entre várias telas (navegação) e como os dados fluem entre as telas (gerenciamento de estado). Ao terminar, você será capaz de montar o esqueleto clássico de um app com várias telas: “lista detalhe edição”.}

\booksection{O panorama do fluxo de dados no gerenciamento de estado}

O @State não é a única ferramenta de estado do SwiftUI. Primeiro, veja este panorama para entender “qual usar e quando”:

\begin{center}
\begin{tabularx}{\linewidth}{|X|X|X|X|}
\hline
\textbf{Ferramenta} & \textbf{Cuida de quê} & \textbf{Quem usa} & \textbf{Cenário típico} \tabularnewline
\hline
@State & Estado temporário privado da tela & A view atual & Contadores, texto de campos \tabularnewline
@Binding & Pegar estado emprestado & A view filha altera a view pai & Componentes personalizados \tabularnewline
@StateObject / @Observable & Lógica compartilhada entre telas & View model (ViewModel) & Carregamento de dados, lógica de negócio \tabularnewline
@EnvironmentObject / .environment & Estado global compartilhado & O app inteiro & Sessão do usuário, configurações de tema \tabularnewline
\hline
\end{tabularx}
\end{center}

\begin{center}\small\color{gray}Panorama das ferramentas de estado do SwiftUI (em 2026, @Observable já é o padrão)\end{center}

\begin{lstlisting}[language=Swift]
// O jeito moderno no iOS 17+: a macro @Observable (recomendada para projetos novos)
import Observation

@Observable
class CounterModel {                 // o view model
    var count = 0
    func increment() { count += 1 }
}

struct CounterView: View {
    @State private var model = CounterModel()   // guarda o modelo

    var body: some View {
        VStack {
            Text("Você tocou \(model.count) vezes")
            Button("Toque aqui") { model.increment() }
        }
    }
}
\end{lstlisting}

Conselho direto para iniciantes: para estado privado da tela, @State; para o que cruza telas ou tem lógica, classes com @Observable. O @StateObject/@ObservableObject dos tutoriais antigos é a geração anterior — saber ler basta; em código novo, siga a orientação acima.

\booksection{NavigationStack: a “pilha” de navegação entre telas}

Desde o iOS 16, a navegação entre várias telas usa o NavigationStack — como uma pilha de pratos: você coloca um em cima e, ao voltar, eles saem um a um:

\begin{lstlisting}[language=Swift]
struct HomeView: View {
    var body: some View {
        NavigationStack {
            List {
                NavigationLink("O livro de João") {
                    DetailView(name: "O livro de João")   // tela de destino
                }
                NavigationLink("O livro de Maria") {
                    DetailView(name: "O livro de Maria")
                }
            }
            .navigationTitle("Lista de livros")   // título grande no topo
        }
    }
}

struct DetailView: View {
    let name: String
    var body: some View {
        Text("Detalhes de \(name)")
            .navigationTitle(name)
            .navigationBarTitleDisplayMode(.inline)
    }
}
\end{lstlisting}

O NavigationLink é o “link de navegação clicável”; o navigationTitle gera automaticamente o título da barra de navegação e o botão de voltar. Para navegar por código (por exemplo, voltar automaticamente depois de tocar em “Salvar”), use @Environment(\textbackslash{}.dismiss) para fechar a tela atual — com essas duas técnicas, 90\% das necessidades de navegação estão cobertas.

\booksection{Janelas modais e abas inferiores: sheet e TabView}

Além de “empurrar telas”, há duas interações clássicas:

\begin{lstlisting}[language=Swift]
// 1. Sobe de baixo em meia tela/tela cheia (o favorito para formulários)
struct HomeView: View {
    @State private var showAdd = false
    var body: some View {
        Button("Nova entrada") { showAdd = true }
            .sheet(isPresented: $showAdd) {
                AddRecordView()     // a tela de entrada que sobe
            }
    }
}

// 2. A barra de abas inferior (o "esqueleto" do app)
struct MainTabView: View {
    var body: some View {
        TabView {
            RecordListView()
                .tabItem { Label("Livro", systemImage: "book") }
            StatsView()
                .tabItem { Label("Estatísticas", systemImage: "chart.pie") }
            SettingsView()
                .tabItem { Label("Ajustes", systemImage: "gear") }
        }
    }
}
\end{lstlisting}

\begin{tcbinfo}{Dica: quando usar cada uma?}
Regra de bolso: no fluxo principal, empurre telas com NavigationStack (lista detalhe); para ações temporárias, use janelas sheet (adicionar, editar, confirmar); para grandes blocos equivalentes, TabView (abas inferiores). A combinação dos três é o esqueleto da grande maioria dos apps.

\end{tcbinfo}

\booksection{List: a “lista nacional” do iOS}

O componente mais usado e mais com “jeitão Apple” do iOS é o List — já vem com divisórias, exclusão por arrasto e destaque ao toque:

\begin{lstlisting}[language=Swift]
struct RecordListView: View {
    @State private var records = ["Almoço 25.5", "Metrô 6.0", "Salário +8000"]

    var body: some View {
        List {
            ForEach(records, id: \.self) { r in
                Text(r)
            }
            .onDelete { indexSet in          // deslize para a esquerda para excluir
                records.remove(atOffsets: indexSet)
            }
        }
    }
}
\end{lstlisting}

Lista + NavigationLink + sheet + TabView — combinando esses quatro elementos, você já consegue montar o esqueleto completo de um app: “lista do diário financeiro abrir o detalhe tocar no + para abrir a janela de lançamento”. É assim que o projeto do próximo capítulo será construído.

\begin{tcbpractice}{Mãos à obra: monte um “esqueleto de três telas”}
\booksection{Ciclo de vida do app: primeiro plano e segundo plano}

Como o ciclo de vida da Activity no Android, um app iOS também passa por estados de “nascimento, primeiro plano, segundo plano e morte”. No SwiftUI, observe pelo valor de ambiente scenePhase:

\begin{lstlisting}[language=Swift]
@main
struct MyApp: App {
    @Environment(\.scenePhase) private var phase

    var body: some Scene {
        WindowGroup {
            ContentView()
        }
        .onChange(of: phase) { _, newPhase in
            switch newPhase {
            case .active:    print("De volta ao primeiro plano, retomando atualização")
            case .inactive:  print("Em transição (chamada / acabou de ir para o fundo)")
            case .background: print("Foi para o fundo, salvando rascunho")
            @unknown default: break
            }
        }
    }
}
\end{lstlisting}

Usos típicos: atualizar dados ao voltar para o primeiro plano e salvar rascunhos ao ir para o segundo plano. O iOS é bem rigoroso com a execução em segundo plano — um app comum é suspenso (congelado) pouco depois de sair da frente, então não conte com processos pesados escondidos lá; se precisar de tarefas agendadas, use as notificações locais do Capítulo 68, e não um “app residente em segundo plano”.

Monte a versão com dados fictícios do esqueleto do diário financeiro: um TabView com três abas (Registros/Estatísticas/Ajustes) a tela de registros usa List para exibir dados fictícios, e tocar numa linha abre o detalhe (NavigationStack) o botão “+” no canto superior direito abre a janela de lançamento em sheet. Sem persistência e sem internet, só para treinar o esqueleto — com ele pronto, encher de dados reais no Capítulo 65 será um pulo.

\end{tcbpractice}

\begin{tcbsummary}{Resumo do capítulo}
\begin{itemize}
  \item As quatro ferramentas de estado: @State (privado da tela), @Binding (emprestado), @Observable (lógica) e .environment (global)
  \item Em projetos novos, escreva view models com o macro @Observable; o @ObservableObject dos tutoriais antigos, basta saber ler
  \item NavigationStack empurra telas; @Environment(\textbackslash{}.dismiss) volta por código
  \item Fluxo principal empurra telas, ações temporárias abrem sheet, blocos equivalentes usam TabView
  \item O List é a “lista nacional”: divisórias e exclusão por arrasto prontas de fábrica
\end{itemize}

\end{tcbsummary}
